\section{Calculate GCD}
\label{sec:n-gcd}
\problemlink{https://atcoder.jp/contests/math-and-algorithm/tasks/math_and_algorithm_o}{AtCoder math\_and\_algorithm 015}
\source{AtCoder Math \& Algorithm 015 (Notion: Number Theory)}{Easy}

\problemstatement{Given $1\le A,B\le10^9$, print $\gcd(A,B)$.}

\begin{patternbox}
The Euclidean algorithm --- the first tool of every number theory problem.
\end{patternbox}

\subsection*{Approach and intuition}

Any common divisor of $a$ and $b$ also divides $a-qb=a\bmod b$, and vice
versa, so $\gcd(a,b)=\gcd(b,a\bmod b)$. The second argument strictly
decreases and reaches $0$, where $\gcd(a,0)=a$. Every two steps the larger
number at least halves, so there are $O(\log\min(a,b))$ steps.

\subsection*{Algorithm}
\begin{algorithmic}[1]
\While{$b\ne0$} \State $(a,b)\gets(b,a\bmod b)$ \EndWhile
\State \Return $a$
\end{algorithmic}

\dryrun{$(33,88)\to(88,33)\to(33,22)\to(22,11)\to(11,0)$: $11$.
$(123,777)\to(777,123)\to(123,39)\to(39,6)\to(6,3)\to(3,0)$: $3$. \checkmark}

\subsection*{C++17 implementation}
\cppfile{number-theory/01_calculate_gcd.cpp}

\begin{complexitybox}
\textbf{Time:} $O(\log\min(A,B))$. \textbf{Space:} $O(1)$.
\end{complexitybox}

\begin{pitfallbox}
\begin{itemize}
  \item In contests just use \texttt{std::gcd} (C++17, \texttt{<numeric>}).
  \item The subtraction-only version (\texttt{a -= b}) is $O(\max)$ ---
        TLE for $(10^9,1)$.
\end{itemize}
\end{pitfallbox}
